%-----------------------------------
% Kapitel 5: Evaluation und zukünftige Technologiepotenziale
%-----------------------------------
\section{Evaluation und zukünftige Technologiepotenziale}
\label{evaluation}

Kapitel~5 tritt in die Evaluationsphase des in Abschnitt~\ref{methodik} dargelegten Vorgehens ein. Abschnitt~\ref{wuerdigung} w\"urdigt das in Kapitel~\ref{sollKonzeption} entwickelte Artefakt kritisch, indem er die erwarteten Wirkungen an den Kennzahlen aus Abschnitt~\ref{kennzahlen} pr\"uft und ihnen die Risiken gegen\"uberstellt. Abschnitt~\ref{ausblick} richtet den Blick auf Entwicklungen, die \"uber den vorliegenden Entwurf hinausweisen. Da keine empirische Erhebung stattfindet, handelt es sich durchgehend um erwartete und nicht um gemessene Wirkungen; die Aussagen betreffen Richtungen, nicht Betr\"age.

\subsection{Kritische Würdigung des Soll-Konzeptes}
\label{wuerdigung}

Die Evaluation im Sinne der Design Science Research Methodology vergleicht das entwickelte Artefakt mit den Zielen, die zu Beginn festgelegt wurden\footcite[Vgl.][S.~56]{Peffers2007}. Venable, Pries-Heje und Baskerville ordnen Evaluationen in der gestaltungsorientierten Forschung nach zwei Dimensionen, n\"amlich dem Zweck, der formativ oder summativ sein kann, und dem Paradigma, das artifiziell oder naturalistisch angelegt ist\footcite[Vgl.][hier S.~77]{Venable2016}. Die nachfolgende W\"urdigung ist in dieser Systematik eine formative Evaluation im artifiziellen Paradigma und findet ex ante statt, also vor einer Umsetzung. Sie pr\"uft das Artefakt argumentativ gegen die Anforderungen und die Befunde der Literatur, nicht im Einsatz bei Anwendern. Das ist keine Notl\"osung, sondern der f\"ur ein noch nicht implementiertes Artefakt vorgesehene Typ; seine Grenze liegt darin, dass er \"uber die Wirkungsrichtung Auskunft gibt und nicht \"uber die Wirkungsh\"ohe. Die Forschungsfrage aus Abschnitt~\ref{zielsetzung} nennt vier Anforderungen, n\"amlich Effizienz, Qualit\"at, Datenschutz und Kontrolle. \autoref{tab:erfuellung} stellt ihnen die Elemente des Soll-Konzeptes gegen\"uber und weist aus, in welchem Ma{\ss}e die Anforderung als erf\"ullt gelten kann.

\begin{table}[htbp]
	\centering
	\caption{Anforderungen der Forschungsfrage und ihre Umsetzung im Soll-Konzept}
	\label{tab:erfuellung}
	\begin{tabularx}{\textwidth}{@{}p{2.1cm}>{\raggedright\arraybackslash}X>{\raggedright\arraybackslash}X@{}}
		\toprule
		Anforderung & Umsetzung im Konzept & Erf\"ullungsgrad \\
		\midrule
		Effizienz & Dialogorientierte Erfassung, maschinelle Kategorisierung und Priorisierung, abrufgest\"utzte L\"osungsvorschl\"age (Abschnitt~\ref{prozessschritte}) & Wirkungsrichtung f\"ur alle vier Kennzahlen hergeleitet, H\"ohe der Wirkung nicht gemessen \\
		Qualit\"at & Konfidenzschwelle vor der \"Ubernahme, Begr\"undungspflicht des Vorschlags, Protokollierung jeder Korrektur (Abschnitte~\ref{kategorisierungSoll} und \ref{fallback}) & konzeptionell verankert, Wirksamkeit von der Kalibrierung des Modells abh\"angig \\
		Datenschutz & Positivliste f\"ur vollautomatisierte Antworten, Ausschluss automatisierter Bearbeitung bei sensiblen Inhalten, Zweckbindung der Korrekturprotokolle (Abschnitte~\ref{voraussetzungen} und \ref{fallback}) & konzeptionell verankert, Umsetzung setzt organisatorische Festlegungen voraus \\
		Kontrolle & Gestaffelte Beteiligung des Menschen, menschliche Entscheidungsknoten, Eskalationsentscheidung beim Agenten (Abschnitte~\ref{hitl} und \ref{entscheidungsknoten}) & im Prozessmodell abgebildet, Wirksamkeit von der tats\"achlichen Pr\"ufung durch die Agenten abh\"angig \\
		\bottomrule
	\end{tabularx}
\end{table}

Die Matrix zeigt ein einheitliches Muster. Alle vier Anforderungen sind im Prozessmodell verankert und nicht als begleitende Ma{\ss}nahmen angef\"ugt; keine von ihnen ist jedoch allein durch die Konzeption eingel\"ost, da ihre Wirksamkeit von Gr\"o{\ss}en abh\"angt, die erst im Betrieb entstehen. Abschnitt~\ref{effizienz} leitet die erwarteten Effizienzwirkungen f\"ur die vier Kennzahlen ab und benennt dabei auch gegenl\"aufige Effekte. Abschnitt~\ref{herausforderungen} behandelt anschlie{\ss}end die Herausforderungen, die dem Konzept Grenzen setzen.

\subsubsection{Erwartete Effizienzsteigerungen anhand der in Abschnitt~3.3 definierten Kennzahlen}
\label{effizienz}

F\"ur die First Contact Resolution Rate wirken zwei Mechanismen in dieselbe Richtung. Die bedeutungsbasierte Suche aus Abschnitt~\ref{loesungsvorschlaege} findet Wissensartikel auch dann, wenn der Anwender eine andere Terminologie verwendet als deren Autoren, womit das in Abschnitt~\ref{unstrukturiert} beschriebene Stichwortproblem entf\"allt und mehr F\"alle bereits im Erstkontakt l\"osbar werden. Hinzu tritt der Self-Service-Pfad, der Standardanliegen ohne Beteiligung eines Agenten abschlie{\ss}t und sie damit definitionsgem\"a{\ss} im ersten Kontakt l\"ost.

Die Kennzahl reagiert auf diese Ver\"anderung allerdings nicht eindeutig. Werden Standardf\"alle automatisiert abgeschlossen, verbleiben im Agenten-Kanal \"uberproportional die schwierigeren Vorg\"ange, deren L\"osung im Erstkontakt seltener gelingt. Die f\"ur diesen Kanal gemessene FCR kann dadurch sinken, obwohl die Leistung des Gesamtprozesses steigt. Eine belastbare Interpretation setzt deshalb voraus, dass die Kanalverschiebung mit ausgewiesen wird, etwa durch getrennte Betrachtung von Self-Service- und Agenten-Pfad neben einem zusammengefassten Wert \"uber beide.

Bei der Mean Time to Resolve greifen vier Mechanismen ineinander. Die manuelle Triage aus Abschnitt~\ref{manuelleKategorisierung} entf\"allt als eigener Zeitblock, sofern der Vorschlag angenommen wird. Die asynchronen R\"uckfrageschleifen verk\"urzen sich, weil die Erfassung fehlende Angaben bereits im Dialog erhebt. Eine treffendere Kategorisierung senkt die Zahl der Fehlzuweisungen und damit der Ping-Pong-Eskalationen. Der Kontexttransfer aus Abschnitt~\ref{kontexttransfer} beschleunigt schlie{\ss}lich die \"Ubergabe an den Second Level, weil dort weniger nachgefragt und weniger doppelt diagnostiziert wird.

Einen empirischen Anhaltspunkt f\"ur die Gr\"o{\ss}enordnung liefert die bereits in Abschnitt~\ref{ausgangssituation} herangezogene Untersuchung, die f\"ur LLM-gest\"utzte Assistenz eine Zunahme der je Stunde gel\"osten F\"alle um 15~Prozent ausweist und den st\"arksten Effekt bei unerfahrenen Mitarbeitenden feststellt\footcite[Vgl.][S.~890]{Brynjolfsson2025}. \"Ubertragen auf den vorliegenden Entwurf legt dieser Befund nahe, dass der Nutzen dort am gr\"o{\ss}ten ausf\"allt, wo die Fluktuation hoch und der Anteil neu eingearbeiteter Agenten entsprechend gro{\ss} ist. Zugleich mahnt die Herkunft des Befundes zur Zur\"uckhaltung, denn er stammt aus einem Kundenservice-Umfeld und nicht aus einem an ITIL ausgerichteten Service Desk.

Die First Level Resolution Rate d\"urfte am deutlichsten zunehmen, da beide Schwachstellenfelder unmittelbar auf sie wirken. Treffendere Kategorien verringern die Zahl der Vorg\"ange, die allein wegen einer Fehlzuordnung den First Level verlassen, und abrufgest\"utzte L\"osungsvorschl\"age erweitern den Kreis der dort abschlie{\ss}bar bearbeitbaren F\"alle. Gegenl\"aufig wirkt allein die in Abschnitt~\ref{fallback} vorgesehene Eskalation bei unzureichender Konfidenz, die Vorg\"ange bewusst fr\"uher weitergibt. Da f\"ur die FLR kein \"offentlich belegter Vergleichswert vorliegt, ist ihre Entwicklung ausschlie{\ss}lich gegen den Ausgangswert der jeweiligen Organisation zu beurteilen.

Beide Wirkungen h\"angen vom Kanalmix ab. Da telefonische Meldungen nach Abschnitt~\ref{prozessarchitektur} weiterhin vom Agenten erfasst werden, f\"allt der Erfassungsgewinn bei hohem Telefonanteil geringer aus. Das Ticketvolumen pro Agent d\"urfte steigen, da Erfassung und Triage weitgehend entfallen und die Recherche k\"urzer ausf\"allt. Die in Abschnitt~\ref{kennzahlen} formulierte Einschr\"ankung gilt hier jedoch versch\"arft. Die Kennzahl ist ohne Betrachtung des Ticketmix wenig aussagekr\"aftig, und genau diesen Mix ver\"andert der Self-Service-Pfad, indem er einfache Vorg\"ange herausl\"ost. Ein gestiegener Wert w\"are daher nicht ohne Weiteres als Produktivit\"atsgewinn zu lesen; ebenso denkbar ist, dass die verbleibenden F\"alle je Vorgang mehr Zeit beanspruchen und die Kennzahl trotz besserer Werkzeuge stagniert.

S\"amtliche Aussagen dieses Abschnitts beschreiben Wirkungsrichtungen, die sich aus der Konstruktion des Prozesses ergeben, und keine gemessenen Effekte. Die in Abschnitt~\ref{kennzahlen} referierten Vergleichswerte taugen als Bezugsrahmen f\"ur eine sp\"atere Einordnung, nicht als Prognose f\"ur eine konkrete Organisation\footcites[Vgl.][o.~S.\isdot]{Rumburg2011}[][o.~S.\isdot]{Rumburg2018}[][o.~S.]{MetricNet2020}. Eine belastbare Quantifizierung setzt eine Pilotierung mit Vorher-Nachher-Messung voraus, in der die vier Kennzahlen vor und nach Aktivierung der KI-Komponenten unter sonst vergleichbaren Bedingungen erhoben werden. Diese Grenze ist der literaturbasierten Anlage der Arbeit geschuldet.

Effizienz ist zudem ein Verh\"altnis von Ertrag und Aufwand, weshalb die bislang betrachtete Ertragsseite allein kein Urteil erlaubt. Auf der Aufwandsseite stehen vier Bl\"ocke. Einmalig f\"allt die Einf\"uhrung an, also die Anbindung an das ITSM-System, die Feinabstimmung des Kategorisierungsmodells auf den vorhandenen Ticketbestand und die Aufbereitung der Wissensdatenbank f\"ur den Abruf. Laufend entstehen Betriebskosten f\"ur die Modellnutzung, die bei einem extern bezogenen Dienst mit der Zahl der Anfragen wachsen und bei lokalem Betrieb als Vorhaltung von Rechenkapazit\"at anfallen. Hinzu tritt der fortlaufende Pflegeaufwand, da der in Abschnitt~\ref{loesungsvorschlaege} beschriebene Abruf nur so gut ist wie der Wissensbestand, auf den er zugreift. Schlie{\ss}lich verursachen die Schulung der Agenten und die in Abschnitt~\ref{voraussetzungen} genannten organisatorischen Festlegungen Aufwand, der nicht technischer Natur ist.

Eine Bezifferung dieser Bl\"ocke ist ohne konkreten Organisationsbezug nicht m\"oglich, da sie vom Ticketvolumen, vom gew\"ahlten Betriebsmodell und vom Zustand der vorhandenen Wissensdatenbank abh\"angt. Festhalten l\"asst sich jedoch die Struktur des Zusammenhangs: Der Einf\"uhrungsaufwand f\"allt unabh\"angig vom Ticketvolumen an, w\"ahrend die Entlastung mit ihm w\"achst. Der Aufbau lohnt sich daher erst oberhalb eines Volumens, das im Einzelfall zu bestimmen ist. Den erwarteten Nutzeffekten stehen \"uberdies Risiken gegen\"uber, die der folgende Abschnitt behandelt.

\subsubsection{Herausforderungen (Halluzinationen von LLMs, Datenschutz, IT-Sicherheit, Akzeptanz)}
\label{herausforderungen}

Sprachmodelle erzeugen Inhalte, die durch die zugrunde liegenden Quellen nicht gedeckt sind. Zhao et al.~unterscheiden intrinsische Halluzinationen, die der bereitgestellten Eingabe widersprechen, von extrinsischen, die sich an ihr nicht \"uberpr\"ufen lassen\footcite[Vgl.][S.~59]{Zhao2023}. F\"ur den Support wiegt die erste Kategorie schwerer, weil ein L\"osungsvorschlag an den abgerufenen Wissensartikeln zu messen ist, die dem Modell als Eingabe vorliegen.

Im Soll-Prozess wirkt sich das an zwei Stellen aus. Ein unzutreffender, aber plausibel formulierter L\"osungsvorschlag kann einen Agenten zu einem Eingriff veranlassen, der die St\"orung versch\"arft. Eine fehlerhafte Self-Service-Antwort erreicht den Anwender ohne Zwischenpr\"ufung. Das Konzept begegnet dem mit der Bindung an abgerufene Wissensartikel, der Ausweisung der Fundstellen, der Pr\"ufung durch den Agenten und der Beschr\"ankung des Self-Service auf eine Positivliste. Diese Ma{\ss}nahmen greifen ineinander, beseitigen das Problem jedoch nicht. Die Quellenbindung senkt die Wahrscheinlichkeit unzutreffender Aussagen, schlie{\ss}t sie aber nicht aus, etwa wenn abgerufene Artikel veraltet sind oder das Modell sie unzutreffend verdichtet. Die menschliche Pr\"ufung wiederum wirkt nur, soweit sie tats\"achlich stattfindet; der in Abschnitt~\ref{fallback} herangezogene Befund zur schlecht kalibrierten Nutzung algorithmischer Vorschl\"age gilt hier unver\"andert.\footcite[Vgl.][S.~1]{GreenChen2019}

Das zweite Feld betrifft den Datenschutz. Anwender liefern im Freitext regelm\"a{\ss}ig Angaben mit, nach denen niemand gefragt hat und deren Sensibilit\"at sich erst bei der Verarbeitung zeigt. Da die Erfassung nach Abschnitt~\ref{ticketerfassung} gerade auf offene Schilderungen setzt, beg\"unstigt der Entwurf dieses Verhalten sogar. Bei externem Modellbetrieb treten Fragen der \"Ubermittlung, der Zweckbindung und der Speicherdauer hinzu, f\"ur die Abschnitt~\ref{voraussetzungen} den rechtlichen Rahmen benannt hat. Offen bleiben mehrere Punkte, die eine Umsetzung zu kl\"aren h\"atte. Das Dialogdesign w\"are so anzulegen, dass es nicht mehr Angaben hervorruft als erforderlich. F\"ur Dialoghistorien und Modellprotokolle w\"aren L\"oschfristen festzulegen. Vor Inbetriebnahme ist eine Datenschutz-Folgenabsch\"atzung durchzuf\"uhren, sofern die Verarbeitung voraussichtlich ein hohes Risiko f\"ur die Rechte und Freiheiten nat\"urlicher Personen zur Folge hat; die Aufsichtsbeh\"orden f\"uhren hierzu Listen der Verarbeitungsvorg\"ange, f\"ur die eine solche Absch\"atzung verpflichtend ist\footcites[Vgl.][Art.~35 Abs.~1]{DSGVO}[][Art.~35 Abs.~4]{DSGVO}. Die Liste der Datenschutzkonferenz f\"ur den nicht-\"offentlichen Bereich f\"uhrt unter Nr.~11 den Einsatz k\"unstlicher Intelligenz zur Verarbeitung personenbezogener Daten auf, soweit er die Interaktion mit den Betroffenen steuert, und nennt als typisches Einsatzfeld ausdr\"ucklich den Kundensupport mittels k\"unstlicher Intelligenz\footcite[Vgl.][Nr.~11, S.~3]{DSK2018}. Die dialogorientierte Erfassung des Soll-Prozesses entspricht dieser Beschreibung, auch wenn die Betroffenen hier Besch\"aftigte und keine Kunden sind; eine Datenschutz-Folgenabsch\"atzung ist deshalb vor Inbetriebnahme durchzuf\"uhren. F\"ur die Korrekturprotokolle ist zus\"atzlich Nr.~8 der Liste einschl\"agig, die umfangreiche Verarbeitungen von Daten \"uber das Verhalten Besch\"aftigter erfasst, sofern diese zur Bewertung der Arbeitst\"atigkeit so eingesetzt werden k\"onnen, dass die Betroffenen erheblich beeintr\"achtigt werden\footcite[Vgl.][Nr.~8, S.~3]{DSK2018}. Ob dieser Tatbestand greift, h\"angt vom Umfang der Protokollierung und von ihrer Verwendbarkeit ab; die in Abschnitt~\ref{voraussetzungen} festgelegte Zweckbindung setzt genau an dieser Stelle an. Diese Punkte sind im vorliegenden Konzept als Anforderung benannt, nicht ausgearbeitet.

Das dritte Feld betrifft die IT-Sicherheit und folgt unmittelbar aus der Architektur des Soll-Prozesses. Das Sprachmodell verarbeitet Freitext, den ein Anwender verfasst, und ruft zus\"atzlich Inhalte aus der Wissensdatenbank ab. Beide Wege sind Einfallstore f\"ur Prompt Injection, also f\"ur Eingaben, die das Verhalten des Modells in unbeabsichtigter Weise ver\"andern. Die OWASP Foundation f\"uhrt diese Schwachstelle als gr\"o{\ss}tes Risiko von Anwendungen mit Sprachmodellen und unterscheidet die direkte Einschleusung \"uber die Nutzereingabe von der indirekten \"uber abgerufene Inhalte; abrufgest\"utzte Verfahren beseitigten das Problem nicht\footcite[Vgl.][o.~S.]{OWASP2025}. F\"ur den vorliegenden Prozess bedeutet dies zweierlei. Eine St\"orungsmeldung k\"onnte Anweisungen enthalten, die auf eine unzutreffend hohe Priorit\"at oder auf die Ausgabe von Inhalten zielen, die dem Melder nicht zustehen. Und ein manipulierter Wissensartikel w\"urde \"uber den Abruf in die Antwort gelangen, ohne dass die Manipulation im Dialog sichtbar w\"urde. Der Entwurf begegnet dem an drei Stellen, die bereits angelegt sind: Die Positivliste begrenzt die vollautomatisierte Beantwortung auf einen definierten Katalog, die menschliche Pr\"ufung an den \"ubrigen Knoten f\"angt auff\"allige Ausgaben ab, und die Konfidenzanzeige macht unsichere Vorschl\"age erkennbar. Hinzutreten m\"ussen eine als unsicher behandelte Trennung von Nutzereingabe und Systemanweisung sowie ein Freigabeprozess f\"ur Wissensartikel, der festlegt, wer Inhalte einstellen darf. Damit ist das Risiko begrenzt, aber nicht beseitigt.

Das vierte Feld betrifft die Akzeptanz, wobei zwei Nutzergruppen zu unterscheiden sind. F\"ur die Agenten bietet die Unified Theory of Acceptance and Use of Technology (UTAUT) einen Analyserahmen; Venkatesh et al.~f\"uhren Leistungserwartung, Aufwandserwartung, sozialen Einfluss und erleichternde Bedingungen als Bestimmungsgr\"o{\ss}en der Nutzungsabsicht an\footcite[Vgl.][hier S.~447]{Venkatesh2003}. Die Leistungserwartung d\"urfte hier ambivalent ausfallen, da ein System, das Kategorien vorschl\"agt und Korrekturen protokolliert, zugleich als Entlastung und als Kontrolle der eigenen Arbeit wahrgenommen werden kann. Hinzu treten zwei gegenl\"aufige Fehlhaltungen. Dietvorst, Simmons und Massey zeigen, dass Personen algorithmische Verfahren nach beobachteten Fehlern mieden, selbst wenn diese dem menschlichen Urteil \"uberlegen blieben\footcite[Vgl.][hier S.~115]{Dietvorst2015}. Parasuraman und Manzey beschreiben umgekehrt die unkritische \"Ubernahme automatisierter Ausgaben infolge nachlassender Aufmerksamkeit\footcite[Vgl.][hier S.~381]{ParasuramanManzey2010}. Beide Haltungen f\"uhren zu unerw\"unschten Ergebnissen, die sich nicht allein technisch auffangen lassen; erforderlich sind Schulung, ein offener Umgang mit Fehlern des Systems und die in Abschnitt~\ref{entscheidungsknoten} vorgesehene Konfidenzanzeige. F\"ur die Anwender stellt sich die Frage anders. Ein Self-Service l\"asst sich als Beschleunigung erleben oder als Abschottung der IT-Organisation, je nachdem, wie leicht ein Mensch erreichbar bleibt. Die in Abschnitt~\ref{ticketerfassung} vorgesehene Ausweichoption ist damit weniger eine Komfortfunktion als eine Bedingung der Akzeptanz.

Die vier Felder begrenzen den erreichbaren Automatisierungsgrad, ohne die hybride Architektur selbst zu ber\"uhren, da sie s\"amtlich an der Verl\"asslichkeit maschineller Ausgaben ansetzen und nicht an der Aufteilung zwischen Maschine und Mensch. Alle vier begrenzen jedoch, wie weit die Automatisierung tragen kann, und zwar unabh\"angig davon, wie leistungsf\"ahig die eingesetzten Modelle sind. Die hybride Architektur erweist sich damit nicht als \"Ubergangsl\"osung auf dem Weg zur Vollautomatisierung, sondern als Konsequenz aus diesen Grenzen. Welche technologischen Entwicklungen diese Grenzen verschieben k\"onnten, behandelt der folgende Abschnitt.

\subsection{Ausblick auf zukünftige Entwicklungen im IT-Support}
\label{ausblick}

Das in Kapitel~\ref{sollKonzeption} entwickelte Konzept beschreibt den Stand dessen, was sich gegenw\"artig verantwortbar umsetzen l\"asst. Zwei Entwicklungslinien k\"onnten die Aufgabenteilung zwischen Mensch und System k\"unftig weiter verschieben. Abschnitt~\ref{kiAgenten} betrachtet den \"Ubergang von assistierenden zu agentischen Systemen, Abschnitt~\ref{predictive} die Verlagerung von der reaktiven zur vorausschauenden St\"orungsbearbeitung.

\subsubsection{Von assistierenden LLMs zu autonomen KI-Agenten}
\label{kiAgenten}

Das Soll-Konzept nutzt Sprachmodelle assistierend. Sie erzeugen Vorschl\"age, w\"ahrend der Prozessrahmen den Ablauf und der Mensch die Entscheidung bestimmt. KI-Agenten gehen dar\"uber hinaus. Wang et al.~beschreiben sie als Systeme, die auf Grundlage eines Sprachmodells Aufgaben in Teilschritte zerlegen, externe Werkzeuge aufrufen, Zwischenergebnisse bewerten und ihr Vorgehen anpassen; als architektonische Bestandteile nennen sie ein Planungsmodul, ein Ged\"achtnis f\"ur zur\"uckliegende Schritte, die Anbindung von Werkzeugen sowie die Ausf\"uhrung selbst\footcite[Vgl.][S.~4]{Wang2024}. Der Unterschied zum vorliegenden Entwurf liegt damit weniger im Modell als in der Frage, wer die Abfolge der Schritte festlegt.

F\"ur den IT-Betrieb liegen erste Untersuchungen vor. Roy et al.~erproben LLM-basierte Agenten f\"ur die Ursachenanalyse von Incidents und lassen diese eigenst\"andig Informationen sammeln und Hypothesen pr\"ufen. Ihre Befunde weisen auf Potenzial hin, benennen zugleich aber Grenzen bei der Zuverl\"assigkeit und der Nachvollziehbarkeit mehrschrittiger autonomer Abl\"aufe\footcite[Vgl.][S.~2]{Roy2024}. Gerade der zweite Punkt wiegt im Support schwer, denn ein Ergebnis, dessen Zustandekommen sich nicht rekonstruieren l\"asst, taugt weder f\"ur die \"Ubergabe an den Second Level noch f\"ur die in Abschnitt~\ref{voraussetzungen} geforderte Nachvollziehbarkeit. Hinzu tritt ein wirtschaftlicher Gesichtspunkt, da ein Agent, der je Vorgang mehrere Modellaufrufe und Werkzeugzugriffe ben\"otigt, deutlich h\"ohere Kosten verursacht als ein einzelner Klassifikationsschritt.

Marktforschungsunternehmen prognostizieren eine raschere Entwicklung. Die bereits in Abschnitt~\ref{ausgangssituation} herangezogene Prognose, wonach agentische KI bis 2029 80~Prozent der g\"angigen Kundenservice-Anfragen ohne menschliches Eingreifen l\"ose, stammt von einem Marktforschungsunternehmen und beschreibt eine Erwartung, keinen empirischen Befund\footcite[Vgl.][o.~S.]{Gartner2025}. Sie taugt als Hinweis auf die Richtung der Aufmerksamkeit, nicht als Planungsgrundlage. Zwischen der in Forschungsarbeiten berichteten Zuverl\"assigkeit und der prognostizierten Autonomie besteht gegenw\"artig eine erhebliche L\"ucke.

F\"ur das vorliegende Konzept folgt daraus kein Umbau, sondern eine Ausbaurichtung. Autonomie lie{\ss}e sich graduell erweitern, indem das System risikoarme L\"osungsschritte nicht nur vorschl\"agt, sondern ausf\"uhrt, etwa das Zur\"ucksetzen eines Kennworts oder die Pr\"ufung einer Berechtigung. Die Entscheidungsknoten aus Abschnitt~\ref{entscheidungsknoten} und die Fallback-Kriterien aus Abschnitt~\ref{fallback} bilden daf\"ur bereits den Steuerungsrahmen, da sie festlegen, was maschinell entschieden werden darf und wann ein Vorgang an den Menschen zur\"uckf\"allt. Autonomie w\"achst in dieser Architektur durch das Verschieben der Schwellen und das Erweitern der Positivliste, nicht durch eine neue Prozessstruktur. Zugleich versch\"arfen sich mit steigender Autonomie die in Abschnitt~\ref{praemissen} begr\"undeten Anforderungen. Je mehr Schritte ein System selbstst\"andig ausf\"uhrt, desto weniger l\"asst sich sein Verhalten aus dem Ergebnis erschlie{\ss}en, weshalb Protokollierung und definierte Eingriffspunkte an Gewicht gewinnen. Die zweite Entwicklungslinie setzt nicht an der Autonomie an, sondern am Zeitpunkt des Eingreifens.

\subsubsection{Proaktives Incident Management (Predictive Support)}
\label{predictive}

Das Soll-Konzept bleibt reaktiv, denn es beginnt mit der Meldung einer St\"orung. Proaktives Incident Management setzt davor an und wertet Betriebsdaten aus Monitoring, Protokolldateien und Ereignismeldungen aus, um Anomalien und Ausfallrisiken zu erkennen, bevor Anwender eine Beeintr\"achtigung bemerken. In der in Abschnitt~\ref{aiopsStand} referierten Taxonomie des Failure Management entspricht dies den Aufgabenfeldern der Pr\"avention und der Erkennung, die dem hier betrachteten Feld der Behebung vorgelagert sind\footcite[Vgl.][S.~7]{Notaro2021}. Der Nutzen ergibt sich daraus, dass eine vor dem Anwenderkontakt behobene St\"orung weder Bearbeitungszeit im Support noch Ausfallzeit beim Anwender verursacht.

Die Verfahren dieses Bereichs beruhen \"uberwiegend auf klassischen Methoden des maschinellen Lernens. Zhang et al.~beschreiben in einer aktuellen \"Ubersicht, wie Sprachmodelle diese erg\"anzen k\"onnten, etwa bei der Interpretation unstrukturierter Protokolldaten und bei der Verkn\"upfung heterogener Datenquellen \"uber Systemgrenzen hinweg\footcite[Vgl.][hier S.~1]{Zhang2025}. Damit r\"uckt dieselbe F\"ahigkeit in den Vordergrund, die auch dem vorliegenden Entwurf zugrunde liegt, n\"amlich die Verarbeitung von Text, der nicht f\"ur die maschinelle Auswertung geschrieben wurde. F\"ur den Support bedeutet dies, dass sich Meldungen der Anwender und Signale der Systeme k\"unftig auf einer gemeinsamen Verarbeitungsgrundlage auswerten lie{\ss}en.

F\"ur das Soll-Konzept erg\"abe sich daraus ein weiterer Eingangskanal. Ein erkanntes Ausfallrisiko k\"onnte ein Ticket erzeugen, bevor Meldungen eintreffen, und den bestehenden Ablauf mit vorbelegter Kategorie und Priorit\"at durchlaufen. Der Self-Service lie{\ss}e sich nutzen, um betroffene Anwender aktiv zu informieren, was das Aufkommen gleichlautender Meldungen und damit auch die Mean Time to Resolve senken d\"urfte. Voraussetzung sind allerdings eine hinreichende Datenqualit\"at und eine Monitoring-Landschaft, die solche Signale \"uberhaupt liefert; beides geht \"uber den Rahmen dieser Arbeit hinaus und w\"are organisationsspezifisch zu pr\"ufen. Zu kl\"aren w\"are \"uberdies der Umgang mit Fehlalarmen, da jedes automatisch erzeugte Ticket ohne tats\"achliche St\"orung genau jene Kapazit\"at bindet, die der Ansatz freisetzen soll.

Beide Entwicklungslinien best\"atigen die Grundentscheidung des Konzepts. Ob die Arbeit die eingangs gestellte Forschungsfrage beantwortet, pr\"uft das folgende Kapitel.
